\noindent\textbf{Context and objective.}
Containers and orchestration platforms require applications to be observed
without modifying their code or relying solely on their logs. eBPF allows
verified programs to be attached to selected points in the Linux kernel,
exposing processes, credentials, files, memory and isolation mechanisms. This
thesis designs and implements a runtime monitoring prototype that selectively
collects host-level events and transfers them to user space with the context
required for analysis.

The reference environment is Rocky Linux 8.10 with kernel 4.18. Compatibility
is verified on the actual kernel because CO-RE and BTF adapt structure access
but cannot provide missing hooks, helpers or attachment mechanisms. The
prototype observes and detects events without changing kernel decisions.

\medskip
\noindent\textbf{Host-level hook implementation.}
The main contribution is the selection and implementation of the host hooks.
The eBPF programs attach to system calls, tracepoints, kernel functions and
Linux Security Module interfaces. A system-call entry preserves the request,
an LSM hook observes security mediation, a tracepoint may confirm a transition,
and a return hook adds the reported result.

Direct hooks produce an event in one invocation. Paired hooks instead retain
entry arguments and complete them at return, using temporary state identified
by thread and event. Per-CPU workspaces keep large records off the eBPF stack.
Arguments are serialised with explicit indexes and bounds, and only the
populated record prefix crosses the perf-event buffer.

Process hooks distinguish a request from the task's actual lifecycle. Fork,
vfork and clone retain the system-call result, while scheduler tracepoints
describe creation, executable replacement and exit. For execution, entry hooks
preserve the path and arguments, \emph{security\_bprm} hooks add the resolved
binary identity, and \emph{sched\_process\_exec} confirms a successful
transition. A separate event records a negative execve or execveat result.

Identity hooks similarly separate request, mediation and installed state. The
setuid and setgid families record identifiers and results;
\emph{security\_task\_fix\_setuid} observes the proposal, whereas
\emph{commit\_creds} compares old and new credentials when they become active.
UIDs, GIDs, the user namespace and capability sets expose the transition
without confusing an intermediate check with its final outcome.

\newpage
\summaryPageHeader

Filesystem hooks distinguish the pathname supplied by the process from the
object resolved by the kernel. The open, chmod and chown system calls describe
request and result, combining several variants in a shared schema.
\emph{security\_file\_open} instead records the resolved path, device, inode and
change time, while retaining the caller pathname when recoverable. Further
hooks observe permissions, ioctl, directory-entry changes and mount topology.

Coverage also includes memory, cross-process transfers, namespaces and cgroups.
The mmap and mprotect system calls preserve requests and results, while security
hooks add the kernel objects involved. \emph{switch\_task\_ns} records only
changed identifiers, and cgroup hooks use local kernel identities.
Kernel-facing surfaces include modules, eBPF operations, kprobes, symbols,
procfs, user-mode helpers and files read by the kernel. These events provide
evidence to detectors without independently proving authorisation or integrity.

\begin{figure}[ht]
    \centering
    \renewcommand{\arraystretch}{1.2}

    \begin{tabular}{ccc}
        \flowbox{Ciclo di vita dei processi\\e transizioni eseguibili} &
        \flowbox{Identita, credenziali\\e controllo dei processi} &
        \flowbox{File e sicurezza\\del filesystem} \\
        \flowbox{Memoria, namespace\\e cgroup} &
        \flowbox{Moduli, eBPF e\\interfacce kernel} &
        \flowbox{Helper comuni\\e filtro UID}
    \end{tabular}

    \vspace{0.55em}
    \textbf{record binario indicizzato e trasporto perf-event}
    \vspace{0.55em}

    \begin{tabular}{ccc}
        \flowbox{Decoder Go\\e registro degli eventi} &
        \flowbox{Policy e detector\\point e collective} &
        \flowbox{Eventi, alert e\\diagnostica separata}
    \end{tabular}

    \caption{Famiglie di hook host-level e percorso comune verso l'analisi in user space.}
    \label{fig:summary-host-hook-families}
\end{figure}


The hooks share initialisation, UID filtering, serialisation and submission.
Their common context preserves process identity, namespaces, timestamps, system
calls and logical events. Explicit CO-RE branches handle differences in the
target kernel, while loops and variable accesses obey verifier-visible bounds.

\medskip
\noindent\textbf{User-space processing.}
The Go runtime selects programs, expands dependencies and decodes records into
typed events through a shared registry. Policies define scope and suppression;
deterministic detectors recognise point conditions or short collective
sequences correlated by process, file or cgroup. Alerts preserve the matching
evidence and may include MITRE ATT\&CK classification metadata.

\medskip
\noindent\textbf{Validation and results.}
The Minikube validation separated the monitor, AES-beta target and client into
three Pods on one node. Three benign runs produced no alerts; three exploit
runs recovered the synthetic flag and generated the expected sequence. This
confirms node-local visibility in the controlled scenario, without implying
general attack coverage or cross-node monitoring.
